\documentclass[10pt,a4paper]{article}

\usepackage[T1]{fontenc}
\usepackage[utf8]{inputenc}
\usepackage{lmodern}
\usepackage{amsmath,amssymb,amsthm}
\usepackage{booktabs,array}
\usepackage{geometry}
\usepackage{xcolor}
\usepackage{enumitem}
\usepackage{microtype}
\usepackage{hyperref}
\usepackage{fancyhdr}

\geometry{margin=1in}
\hypersetup{
  colorlinks=true,
  linkcolor=blue!55!black,
  citecolor=blue!55!black,
  urlcolor=blue!55!black,
  pdftitle={A Bitter End for Bittersweet: Practical Level-I Secret-Key Recovery from Accepted Carry Comparisons},
  pdfauthor={Martin Feussner}
}
\setlist{nosep,leftmargin=*}
\setlength{\parindent}{1em}
\setlength{\parskip}{0.25em}
\pagestyle{fancy}
\fancyhf{}
\fancyhead[L]{Secret-key recovery against Bittersweet Level I}
\fancyhead[R]{8 October 2026}
\fancyfoot[C]{\thepage}
\renewcommand{\headrulewidth}{0.3pt}

\newtheorem{observation}{Observation}
\newtheorem{theorem}{Theorem}
\newcommand{\Z}{\mathbb{Z}}
\newcommand{\roundp}[2]{\left\lfloor #1 \right\rfloor_{2^{#2}}}

\title{\bfseries A Bitter End for Bittersweet\\
\large Practical Level-I Secret-Key Recovery from Accepted Carry Comparisons}
\author{Martin Feussner\\
\small Selmer Center, University of Bergen\\
\small \texttt{martin.feussner@uib.no}}
\date{8 October 2026}

\begin{document}
\maketitle

\begin{center}
\fbox{\parbox{0.92\linewidth}{\small\textbf{AI provenance and review status.}
Target screening, attack-history searches, attack derivation,
implementation, experiments, adversarial review, and the initial manuscript
were produced in an AI-run audit by OpenAI Codex and purpose-built specialist
AI agents.  OpenAI Codex generated and selected the paper title, including its
scheme-specific wordplay.  A separately instructed fresh-context agent executed the frozen
16-signature recovery once on a new key, and another agent independently
replayed its public-only result.  The technical, adversarial,
independent-reproduction, final-history, and publication reviews passed; this
report is classified \texttt{VERIFIED\_ATTACK}.}}
\end{center}
\vspace{0.7em}

\begin{abstract}
Bittersweet is an LWR-based MPC-in-the-head signature whose optimized
responses disclose an opened key-share aggregate and a one-bit reconstruction
error.  At the published Level-I $d=32$ parameters, those values reveal an
oriented strict comparison with a fixed 90-bit buffer of every secret LWR
product.  We collect 16 accepted signatures, giving 416 comparisons for each
of 747 public rows, and intersect them into modular product intervals.  A
weighted modular closest-vector instance on the 160 narrowest intervals is
solved with plain LLL and Babai nearest plane.

In a precommitted held-out execution and a separate fresh-context reproduction,
each public-only candidate satisfied all 747 intervals and rounded public
outputs before private truth was opened.  Hash-pinned comparators reported an
exact match in all eleven 96-bit key words.  Each recovered candidate produced
a signature on a distinct seventeenth message.  The ordinary strict verifier
accepted all 26 repetitions and all 747 rows per repetition, while ten final
mutations rejected.  The reproduction used exactly one new key and one attack
trial.  Its public pipeline took 313.95 seconds and at most 236,760 KiB RSS on
one worker.

The result concerns ePrint 2026/397 version~1's Level-I $d=32$ row under a
frozen canonical realization of the paper's arithmetic and protocol structure.
The paper publishes no author implementation or normative byte format.
Fresh-context reproduction and the dated final history search have passed.
The result is classified \texttt{VERIFIED\_ATTACK} after technical,
adversarial, independent-reproduction, history, and publication review.
\end{abstract}

\noindent\fcolorbox{black!45}{black!4}{\parbox{0.94\textwidth}{
\textbf{Status: VERIFIED\_ATTACK.}  Run-003 and independent run-005 pass the
practical full-parameter recovery and fresh-signing checks. The final dated
attack-history recheck and publication review also pass.}}

\section{Introduction}

Bittersweet~\cite{bittersweet} builds an MPC-in-the-head signature from a
power-of-two Learning With Rounding (LWR) function.  Its public key contains a
three-bit rounding of a matrix--secret product.  To reconcile additive key
sharing with rounding, each optimized response includes a one-bit correction
for every public row.  Signing aborts when a 90-bit slice of the opened-share
product equals the corresponding slice of the full secret product.

The target paper explicitly discusses carry leakage and observes that an
accepted transcript excludes buffer values.  It heuristically posits that the
resulting leakage remains harmless when the signing-query count is very small
relative to $2^{90}$.  We show that far fewer queries suffice at the claimed
Level-I parameters.  The correction bit gives the direction of every excluded
threshold.  Repeating the comparison 416 times per row yields intervals narrow
enough to recover the complete key with conventional lattice tools.

Our concrete contributions are:
\begin{itemize}
  \item an exact derivation from each accepted optimized response to an
  oriented strict comparison with the fixed secret-product buffer;
  \item an interval accumulator using all 16 signatures and 26 repetitions,
  for 310,752 public comparisons;
  \item a practical weighted modular CVP instantiation using 160 rows, plain
  LLL~\cite{lll}, and Babai nearest plane~\cite{babai};
  \item a precommitted held-out execution in which the candidate is frozen and
  independently replayed before truth access;
  \item a fresh-context one-key, one-trial reproduction with independent
  public-only result replay; and
  \item a complete fresh-message signature accepted by the ordinary strict
  verifier, together with mutation controls.
\end{itemize}

No originality is claimed for LLL, Babai, generic modular-CVP embeddings,
power-of-two LWR cryptanalysis, or the general existence of carry leakage.
We describe the modular-interval and accepted-transcript steps as conventional
or target-disclosed components and make no separate novelty claim for them.

\paragraph{Attack-history status.}
No prior public attack against Bittersweet Signatures was located in the
searches performed as of 2026-10-07.  The final post-reproduction search
identified ePrint 2026/250~\cite{powertwolwr} as generic prior
power-of-two-LWR technique, not an attack on exact Bittersweet.  This dated
finding is not an absolute first-attack claim.

\section{Target and paper-level model}

The exact target is IACR ePrint 2026/397 version~1, whose pinned PDF has
SHA-256
\begin{center}
\small\ttfamily eefe27b903cb61fa50fecc33eef3fa99b121bece1c9831a4b3cab7bdebba5dbb.
\end{center}
We execute Table~2's Level-I $d=32$ row:
\[
(\lambda,q,p',\alpha,p,n,m,d,t)
=(128,96,3,90,93,11,747,32,26),
\]
where $p=p'+\alpha$.  The secret is $k\in\Z_{2^{96}}^{11}$.  A public input is
expanded to $X\in\Z_{2^{96}}^{747\times 11}$, and the public rounded output is
\[
                 y=\roundp{Xk}{3}.
\]

In each repetition the signer shares $k$ among 32 parties, commits to the
shares, computes local LWR outputs, and opens all but one share selected by the
Fiat--Shamir challenge.  The optimized signature transmits GGM material or a
correction key for the opened shares, the hidden commitment, and a one-bit
reconstruction-error vector.  The verifier reconstructs the opened aggregate,
uses the bit to rebuild the hidden output, rehashes the complete first message,
and applies Figure~3's modular predicate.

The ePrint supplies no public implementation, normative wire encoding,
domain-separation byte strings, XOF-to-matrix convention, or test vectors.  Our
frozen model uses injective encodings and conventional SHAKE/SHA3 components,
implements both optimized GGM opening cases, checks packed-bit padding, and
preserves the paper's exact arithmetic and Fiat--Shamir bindings.  Our claim is
about that mathematical design.  We do not claim byte compatibility with an
unavailable implementation.

\section{Accepted transcripts reveal strict comparisons}

Fix a public row $X_r$ and write its secret product as
\begin{align*}
u_r&=\langle X_r,k\rangle\bmod 2^{96},\\
u_r&=y_r2^{93}+R_r,\\
R_r&=8B_r+\ell_r,
\end{align*}
where $0\leq B_r<2^{90}$ and $0\leq\ell_r<8$.  The three-bit $y_r$ is public.
The 90-bit buffer $B_r$ is fixed by the public row and signing key, and is
therefore reused across all messages, salts, signatures, and repetitions.

For one response, let $K_{\mathrm{open}}$ be the verifier-reconstructed sum of
the 31 opened key shares.  Define
\[
S_r=\langle X_r,K_{\mathrm{open}}\rangle\bmod 2^{96},\qquad
A_r=(S_r\mathbin{\gg}3)\bmod 2^{90}.
\]
The signing abort rejects equality $A_r=B_r$.  The one-bit reconstruction error
$e_r$, interpreted through Figure~3's modular membership rule, supplies the
orientation:
\[
\begin{array}{lll}
e_r=0 &\Longrightarrow A_r<B_r &\Longrightarrow B_r\geq A_r+1,\\
e_r=1 &\Longrightarrow A_r>B_r &\Longrightarrow B_r\leq A_r-1.
\end{array}
\]

\begin{observation}
Every accepted repetition gives one exact strict comparison per public row
with a row-specific fixed buffer.  It does not reveal 90 independent bits, and
the 747 comparisons in one repetition remain coupled through the same opened
key aggregate.
\end{observation}

The derivation uses the stricter formal Figure~3 rule.  Section~7.1's prose
inequality is ambiguous if read as an ordinary signed comparison; our attack
does not rely on that broader reading.

\section{Accumulating modular product intervals}

Initialize $L_r=0$ and $U_r=2^{90}-1$.  For every signature and repetition,
update
\[
\begin{split}
e_r=0 &: \quad L_r\leftarrow\max(L_r,A_r+1),\\
e_r=1 &: \quad U_r\leftarrow\min(U_r,A_r-1).
\end{split}
\]
Sixteen signatures give $16\cdot26=416$ comparisons per row and
$416\cdot747=310{,}752$ comparisons overall.  The final buffer bracket implies
\[
y_r2^{93}+8L_r
\leq \langle X_r,k\rangle\bmod 2^{96}
\leq y_r2^{93}+8(U_r+1)-1.                 \tag{1}
\]
The last seven values in the upper endpoint account for the unknown low three
bits.

Run-003 and run-005 each began from 16 complete serialized optimized signatures
on distinct messages.  The ordinary strict verifier accepted all of them.
Independent replay reconstructed every opened aggregate and packed error
vector and reproduced the 747 intervals.  Run-003 recorded 171/171 applicable
query-signature controls; run-005 separately recorded 169/169.  No secret label
or private truth is needed to compute Equation~(1).

\section{Weighted lattice recovery}

For each row there exists an integer lift $z_r$ such that
\[
\langle X_r,k\rangle-2^{96}z_r\in[\operatorname{lower}_r,
\operatorname{upper}_r].                                  \tag{2}
\]
Let $c_r$ and $w_r$ be the midpoint and width of this interval.  The frozen
procedure selects the 160 narrowest rows and assigns
\[
\omega_r=\max\left(1,
\operatorname{round}\left(\frac{\max_s w_s}{w_r}\right)\right).
\]
It embeds the weighted modular equations and eleven systematic key
coordinates as a closest-vector problem centered on the interval midpoints.
Plain LLL with $\delta=0.99$ and $\eta=0.501$ reduces the basis, after which
Babai nearest plane produces a candidate.  Every candidate is checked against
all 747 intervals and all 747 public rounded outputs; truth is not part of the
stop rule.

The final schedule was fixed before held-out entropy:
\begin{center}
\begin{tabular}{@{}ll@{}}
\toprule
Accepted signatures & 16 \\
Selected rows & 160 narrowest \\
Block schedule & $[0]$ (plain LLL, then Babai) \\
Configured loops & 2 \\
fpylll seed & 20260412 \\
Stop rule & first candidate satisfying every public constraint \\
\bottomrule
\end{tabular}
\end{center}

Run-003 made one block-zero attempt.  The complete solver took 45.674 seconds
with 235,768 KiB peak RSS.  Before truth access, the candidate fit 160/160 selected intervals,
747/747 total intervals, and 747/747 public rounded outputs.

Run-005 independently executed the same family-B schedule exactly once, with
one public-pipeline invocation, one block-zero attempt, and no tuning or retry.
Its candidate also fit 160/160 selected intervals, 747/747 total intervals,
and 747/747 rounded outputs.  A separate public-only reviewer rebuilt all
310,752 comparisons and replayed the solver and candidate constraints.

\section{Final executions and fresh signing}

The target, sources, runtime, query rule, solver schedule, resource limits,
public/private boundary, truth comparator, message rule, signer, and verifier
were frozen before the run-003 held-out entropy existed.  A new key and 16
complete signatures were generated.  The public pipeline verified and
aggregated them, ran the one configured recovery attempt, and serialized a
public candidate freeze.  Independent public-only and adversarial reviewers
replayed that chain before authorizing the truth comparison.

The hash-pinned comparator reported:
\begin{itemize}
  \item exact equality in all eleven 96-bit words;
  \item identical candidate and truth key digests;
  \item identical 747-product digests;
  \item 747/747 candidate and truth interval checks;
  \item 747/747 public-output checks; and
  \item exact public-key reconstruction.
\end{itemize}
It emitted no raw key words.  The recovered/truth key digest is
\begin{center}
\small\ttfamily 8b7136532fc8be787564db67d2bbdf8e027dba822cfc0412377d96f92df426be.
\end{center}

The fresh-message derivation rule was also frozen before held-out entropy.  It
produced a 140-byte seventeenth message, SHA-256
\begin{center}
\small\ttfamily 22fd90f30a746ab8dbffa4041b2ceba0e2c2988241d624fcc664bbd53ae09f2e,
\end{center}
distinct from every oracle query.  The recovered public candidate produced a
complete 26-response signature in one attempt with zero aborts.  The ordinary
strict verifier accepted all repetitions and all 747 rows per repetition.

Ten applicable mutations rejected: challenge index, message, public matrix,
public rounded output, correction key, error bit, GGM path, hidden commitment,
salt, and nonzero packed-error padding.  A strengthened changed-message test
updated its noncryptographic envelope metadata and still rejected at the
recomputed Fiat--Shamir first-message hash.  An independent replay also bound
the mathematical signature body directly to the original public-key artifact.

\begin{table}[t]
\centering
\caption{Run-003 resource measurements; all stages used one worker.}
\begin{tabular}{@{}lrr@{}}
\toprule
Stage & Wall time & Peak RSS \\
\midrule
Public verification, aggregation, recovery, freeze & 320.47 s & 236,792 KiB \\
Truth comparison & 0.20 s & 26,800 KiB \\
Recovered-key fresh signing & 2.25 s & 28,056 KiB \\
Fresh verification and controls & 15.04 s & 29,696 KiB \\
\bottomrule
\end{tabular}
\end{table}

The sealed public resource summary supports the 320.47-second public pipeline
and the 45.674-second complete solver within it.  It does not bind a separate
key-and-signature-generation duration or a complete run-003 stage sum.  Review
and documentation time is excluded.

\subsection{Independent run-005}

The first reproduction campaign, run-004, failed before entropy: an inherited
standard-input version probe entered an interactive Singular prompt.  It ended
after 32.19 seconds at 228,420 KiB RSS with zero keys, signing queries,
verifier calls, solver trials, truth comparisons, and fresh signatures.  The
fix closed standard input, captured output, and imposed a finite timeout.  An
initial run-005 source review then VETOed a mismatch between the declared
first-nonempty-line parser and the implementation.  Rev2 fixed the parser and
passed before entropy.

Run-005 generated one new OS-random key and 16 complete signatures on distinct
messages with zero oracle aborts.  Its single attack trial produced 747/747
interval and rounded-output checks, followed by 11/11 exact words and 747/747
truth intervals through the sealed comparator.  It signed one distinct
140-byte message, whose SHA-256 is
\begin{center}
\small\ttfamily 00e047a068069375376ff8b812ea98098f89ec9a8f0ee7937e818b4ba14ab9d4.
\end{center}
The signature was produced in one attempt with zero aborts.  The strict
verifier accepted 26/26 repetitions and all 747 rows, and ten final controls
rejected.  A separate public-only result reviewer replayed the comparison,
solver, candidate, and fresh-verifier constraints without opening private
state.

\begin{table}[t]
\centering
\caption{Run-005 resource measurements; every stage used one worker.}
\begin{tabular}{@{}lrr@{}}
\toprule
Stage & Wall time & Peak RSS \\
\midrule
Main precommit & 2.60 s & 230,304 KiB \\
Seed/commitment generation & 2.82 s & 230,496 KiB \\
Fresh key and 16 signatures & 42.80 s & 234,948 KiB \\
Public verification, aggregation, recovery, freeze & 313.95 s & 236,760 KiB \\
Truth comparison & 0.20 s & 26,840 KiB \\
Recovered-key fresh signing & 2.25 s & 28,024 KiB \\
Fresh verification and controls & 15.20 s & 29,676 KiB \\
\bottomrule
\end{tabular}
\end{table}

The seven-stage sum is 379.82 seconds under the 4 GiB cap.

\section{Trial accounting}

The empirical record contains heterogeneous, adaptive, and repeated-key
procedures and cannot be reduced to one binomial success estimate.

\begin{description}[leftmargin=0.25in,style=nextline,itemsep=2pt]
\item[Original projected development.]
One tuned training key has two documented successful configurations, S4/r120
and an independently frozen S8 configuration.  Earlier tuning attempts are not
exhaustively enumerated; this is adaptive development, not campaign evidence.
\item[Fresh adaptive S8 development.]
One new development key had two configurations: r112 failed, then an adaptive
r160 configuration recovered the key and enabled fresh signing.
\item[Run-002 infrastructure and quarantine.]
The first driver stopped at a Sage import before LLL/Babai, adding zero solver
trials.  A separate accidental cross-agent retry exited nonzero with no
candidate and is quarantined as a non-result.
\item[Untouched run-002.]
One held-out key received one frozen S8 trial and no recovery: 0/160 selected
intervals, 1/747 total intervals, and 214/747 public outputs fit.  This negative
is immutable.
\item[Robust Family B development.]
Six shared development keys received six trials and six recoveries, four from
projected marginals and two from complete signatures.  This selected S16.
\item[Robust Family A development.]
The same six keys received six separate trials and six recoveries; two needed
BKZ-24.  Family A ran after Family-B truth opened and is not blind evidence.
\item[Robust environment stops and unused families.]
Two Family-B launches stopped before solving because the NGCC environment was
inactive; their reruns are already counted above.  Families C and D were not
run.  The robust subtotal is six unique keys, twelve substantive A/B trials,
twelve development successes, and two pre-solver environment stops.
\item[Final campaign.]
Run-003 used one untouched key, one fixed S16 trial, and one recovery with fresh
signing.  Run-004 was a zero-key, zero-trial terminal preseed infrastructure
attempt.  The initial run-005 package was VETOed preseed; its rev2 execution
used one OS-random key, one fixed S16 trial, and one recovery with fresh
signing.  Thus the final campaign S16 count is 2/2, with no population-rate
claim.
\item[Public-package validation.]
One earlier full harness generated a local key and signatures but stopped in
controls before any solver attempt, so it contributes zero cryptanalytic
trials.  Five corrected full invocations then used five new local keys and
five fixed trials; all five recovered the generated key and completed fresh
signing without a failed cryptanalytic trial, solver retry, or cryptanalytic
retuning.  They shared the fixed cryptanalytic configuration while
noncryptanalytic harness and reporting code changed between validation
snapshots.  These add zero campaign trials or successes.  The early-stop row
is conservative because no failed result or time log survived.  Quick,
all-row, and determinism checks are zero-key public replays.
\end{description}

The demonstrated campaign query cost remains 16 signatures per attack
execution.  Development, campaign, infrastructure, and package-validation
rows remain separate.

\section{Scope and attribution}

The attack is design-level: it follows from the intended opened-share data,
abort condition, reconstruction-error bit, and formal modular verifier.  It
does not use a parser defect, malformed object, side channel, fault, secret
label, or attack-specific acceptance rule.  Only the Level-I $d=32$ row is
claimed, and the demonstrated procedure requires 16 accepted signatures.

Bittersweet itself identifies the carry-leakage surface.  The lattice recovery
uses established methods.  ePrint 2026/250 by Baudrin, Heim Boissier, and
Standaert is relevant generic power-of-two-LWR prior work, not an exact
Bittersweet attack.  We do not market modular-interval recovery, LLL, Babai,
accepted-transcript leakage, order statistics, or the component methods as
new.

\section{Reproducibility}

The required flat interface provides a quick public replay in
\texttt{reproduce.py}, an independent replay of the sole affected row in
\texttt{reproduce-all-rows.py}, and an optional full mode.  The full mode
generates a key locally, collects 16 complete
signatures, verifies and projects them, recovers and freezes a candidate,
compares only after the freeze, signs a distinct message, invokes the ordinary
verifier, and runs negative controls.  Dynamic results record OS, architecture,
actual interpreter and argument vector, the logical command, dependency
versions, runtime, peak RSS, local trials and recoveries, campaign-trial
increment, and public replay counts.  The quick public-only mode
replays a sanitized interval instance and executes recovery rather than loading
a stored key.

The run-005 campaign record supports AMD EPYC, Python 3.12.14, Singular 4.4.1,
one CPU, and a 4 GiB cap.  The sanitized public runner implements the needed
inverse over $\Z/2^{96}\Z$ directly and does not use SageMath or Singular.  Its
hash lock targets CPython 3.12 on Linux x86-64 with glibc at least 2.17 and
pins official binary wheels for fpylll 0.6.4 and cysignals 1.12.6; source
builds are unverified.  The locked install, imports, and an LLL smoke test
passed in a fresh CPython 3.12.3 virtual environment.  Measured reproducer runs
used CPython 3.12.14.
The exact-final-source full test completed one fresh-key execution in 184.47
seconds internal wall time at 83,336 KiB peak RSS on one worker, with all
recovery and fresh-signing checks passing, 25/25 applicable query controls and
10/10 applicable fresh controls rejecting, and no raw candidate words emitted.

\section{AI disclosure}

Target discovery, history searching, scheme reconstruction,
cryptanalysis, implementation, experiment execution, adversarial review, and
writing were AI-assisted.  Separate AI agents reviewed run-003.  Another AI
agent performed the one-trial fresh-context reproduction, and a further agent
replayed its public-only result.  A clean-context AI Publication Reviewer
conditionally passed the frozen package. The result is classified
\texttt{VERIFIED\_ATTACK} on the basis of the recorded technical, adversarial,
independent-reproduction, history, and publication reviews.

\section*{Acknowledgements}

The computations were performed on the Norwegian Research and Education Cloud (NREC), using resources provided by the University of Bergen and the University of Oslo.

\begin{thebibliography}{9}

\bibitem{bittersweet}
B. Balon, G. Brian, S. Faust, C. Hazay, E. Micheli, and F.-X. Standaert,
``Bittersweet Signatures: Bringing LWR to a Picnic for Hardware-Friendly
MPC-in-the-Head,'' IACR Cryptology ePrint Archive, Report 2026/397, version 1,
2026.  Also in \emph{Security and Cryptography for Networks (SCN 2026)},
LNCS 16810, pp.~96--120, 2026.
\url{https://eprint.iacr.org/2026/397};
\url{https://doi.org/10.1007/978-3-032-36264-3_6}.

\bibitem{lll}
A. K. Lenstra, H. W. Lenstra, Jr., and L. Lov\'asz,
``Factoring polynomials with rational coefficients,''
\emph{Mathematische Annalen}, vol.~261, pp.~515--534, 1982.

\bibitem{babai}
L. Babai,
``On Lov\'asz' lattice reduction and the nearest lattice point problem,''
\emph{Combinatorica}, vol.~6, pp.~1--13, 1986.

\bibitem{powertwolwr}
J. Baudrin, R. Heim Boissier, and F.-X. Standaert,
``On the Concrete Hardness of LWR with a Power of Two Modulus,''
IACR Cryptology ePrint Archive, Report 2026/250, 2026.
\url{https://eprint.iacr.org/2026/250}.

\end{thebibliography}

\end{document}
